\documentclass{article}

\usepackage{amsmath,amssymb}
\usepackage{xurl}
\usepackage[hidelinks]{hyperref}

% Shared commands for notes in docs/paper-gaps/.

\DeclareUnicodeCharacter{2080}{\ifmmode{}_{0}\else\textsubscript{0}\fi}
\DeclareUnicodeCharacter{2081}{\ifmmode{}_{1}\else\textsubscript{1}\fi}
\DeclareUnicodeCharacter{2082}{\ifmmode{}_{2}\else\textsubscript{2}\fi}
\DeclareUnicodeCharacter{2099}{\ifmmode{}_{n}\else\textsubscript{n}\fi}

\newcommand{\Lean}{\textsc{Lean}}
\ExplSyntaxOn
\NewDocumentCommand{\leanid}{m}{
  \ifmmode
    \textsf{\detokenize{#1}}
  \else
    \group_begin:
    \urlstyle{sf}
    \tl_set:Nn \l_tmpa_tl {#1}
    \tl_replace_all:Nnn \l_tmpa_tl {\_} {_}
    \exp_args:NV \path \l_tmpa_tl
    \group_end:
  \fi
}
\ExplSyntaxOff
\providecommand{\tr}{\operatorname{tr}}
\newcommand{\tnleanIssueBase}{https://github.com/LionSR/TNLean/issues/}
\newcommand{\tnleanPrBase}{https://github.com/LionSR/TNLean/pull/}
\newcommand{\tnleanDocBase}{https://sirui-lu.com/TNLean/docs/}
\newcommand{\ghissue}[1]{\href{\tnleanIssueBase#1}{GitHub issue \##1}}
\newcommand{\ghpr}[1]{\href{\tnleanPrBase#1}{PR \##1}}
\newcommand{\doclink}[1]{\href{\tnleanDocBase#1.html}{\path{#1}}}

% Dirac notation and the identity operator, as in blueprint/src/macros/common.tex.
% \providecommand keeps note-local definitions authoritative.
\usepackage{dsfont}
\providecommand{\ket}[1]{|#1\rangle}
\providecommand{\bra}[1]{\langle #1|}
\providecommand{\braket}[2]{\langle #1 | #2 \rangle}
\providecommand{\ketbra}[2]{|#1\rangle\!\langle #2|}
\providecommand{\Id}{\mathds{1}}
\providecommand{\one}{\mathds{1}}
\providecommand{\C}{\mathbb{C}}
\providecommand{\MN}[1]{M_{#1}(\C)}
\providecommand{\MD}{M_D(\C)}

% Verdict marker: \gapnote{<kind>}{<status>}. Kind is the policy
% classification (clarification, local-correction, scope-restriction,
% unfaithful, false-source, open-gap); status is open, wip (elimination
% underway), resolved, or historical. Machine-read by the site index;
% prints a verdict line.
\newcommand{\gapnote}[2]{\par\noindent\textbf{Verdict:} #1 (#2).\par}


\title{Pre-Bialgebras of Matrix Product Operator Algebras: Existence of the Counit of the Fusion Coproduct}
\author{}
\date{2026-11-03}

\begin{document}
\maketitle
\gapnote{scope-restriction}{open}

\section{The source definition}

Ref.~\cite{Molnar2022MPOAlgebrasI} defines a coalgebra as a triple
$(\mathcal C,\Delta,\epsilon)$, where $\mathcal C$ is a finite-dimensional
complex vector space, $\Delta\colon\mathcal C\to\mathcal C\otimes\mathcal C$
is coassociative,
\begin{align}
    (\Delta\otimes\Id)\circ\Delta
        &=(\Id\otimes\Delta)\circ\Delta,
    \label{eq:mrgscp22_coassoc}
\end{align}
and $\epsilon\in\mathcal C^*$ is a counit,
\begin{align}
    (\epsilon\otimes\Id)\circ\Delta
        &=(\Id\otimes\epsilon)\circ\Delta
         =\Id
    \label{eq:mrgscp22_counit}
\end{align}
(Definition~3.1).  A pre-bialgebra is an algebra $\mathcal A$ that is also
a coalgebra such that the coproduct is multiplicative,
\begin{align}
    \Delta(xy)
        &=\Delta(x)\Delta(y)
        \qquad (x,y\in\mathcal A)
    \label{eq:mrgscp22_mult}
\end{align}
(Definition~4.1).\footnote{Source coordinates in \path{mpo.tex} from
    \href{https://arxiv.org/abs/2204.05940}{arXiv:2204.05940}:
    lines~866--875 for Definition~3.1 and lines~1409--1414 for
    Definition~4.1.}
No relation between the counit and the unit or the product is imposed:
the section introduces pre-bialgebras as ``bialgebras without any property
imposed on the unit and counit'', and $\Delta(1)\neq1\otimes1$ is allowed,
in which case the image of $\Delta(1)$ in a tensor product of two
representations is a projector rather than the identity.\footnote{Lines~1407, 1421 and~1708
    (equation \texttt{eq:delta\_epsilon}).}
The counit of~(\ref{eq:mrgscp22_counit}) is nevertheless part of the
definition, and it is used: the unit of the dual pre-bialgebra
$\mathcal A^*$ is the counit of $\mathcal A$ (line~1419), and the tensor
product of two representations of $\mathcal A^*$ is the restriction to the
range of $\Delta(\epsilon)$ (line~1640).

\section{What is formalized}

The formal pre-bialgebra structure on an algebra $\mathcal A$ over a
commutative semiring is the source's: a coalgebra structure, with a
coproduct $\Delta$ satisfying~(\ref{eq:mrgscp22_coassoc}) and a counit
$\epsilon$ satisfying~(\ref{eq:mrgscp22_counit}), whose coproduct is
multiplicative~(\ref{eq:mrgscp22_mult}).  Neither $\Delta(1)=1\otimes1$ nor
any compatibility of $\epsilon$ with the unit or the product is required.
The source works over $\C$ in finite dimension; the formal structure allows
any commutative semiring and any module.

The fusion coproduct of a complete zipper fusion family is defined on the
boundary algebra $\mathcal B=\bigoplus_cM_{D_c}(\C)$ from the fusion
tensors.  It is multiplicative and coassociative, $\Delta(1)$ is the support
projector of the fusion tensors, and on the elements
$T^{ik}=(T_c^{ik})_c$ formed by one letter of every block it is
\begin{align}
    \Delta(T^{ik})
        &=\sum_jT^{ij}\otimes T^{jk}.
    \label{eq:mrgscp22_letter}
\end{align}
These elements span $\mathcal B$.  Consequently, if
$\epsilon\colon\mathcal B\to\C$ is a linear functional with
\begin{align}
    \epsilon(T^{ik})
        &=\delta_{ik},
    \label{eq:mrgscp22_counit_letters}
\end{align}
then~(\ref{eq:mrgscp22_letter}) gives
$(\epsilon\otimes\Id)\Delta(T^{ik})=T^{ik}=(\Id\otimes\epsilon)\Delta(T^{ik})$,
so $\epsilon$ satisfies~(\ref{eq:mrgscp22_counit}) and
$(\mathcal B,\Delta,\epsilon)$ is a pre-bialgebra in the sense of
Definition~4.1.  Such a functional is unique, since the elements $T^{ik}$
span.  It exists when one block $c_0$ is the identity operator,
$T_{c_0}^{ik}=\delta_{ik}\Id$: the normalized trace
$X\mapsto D_{c_0}^{-1}\tr X_{c_0}$ then
satisfies~(\ref{eq:mrgscp22_counit_letters}).\footnote{Formal declarations:
    \leanid{PreBialgebra} in \doclink{TNLean/Algebra/PreBialgebra};
    \leanid{MPOTensor.CompleteZipperFusionFamily.rTensor_counit_comp_comul},
    \leanid{MPOTensor.CompleteZipperFusionFamily.lTensor_counit_comp_comul},
    \leanid{MPOTensor.CompleteZipperFusionFamily.counit_unique},
    \leanid{MPOTensor.CompleteZipperFusionFamily.identityBlockCounit_letter}
    and \leanid{MPOTensor.CompleteZipperFusionFamily.toPreBialgebra} in
    \doclink{TNLean/MPS/MPDO/CompleteZipperFusionCoassoc}.  The blueprint
    entries are \path{def:mpoalg_pre_bialgebra} and
    \path{thm:mpoalg_fusion_coproduct_coassoc} in
    \path{blueprint/src/chapter/ch30_mpo_algebras.tex}.}

The functional~(\ref{eq:mrgscp22_counit_letters}) is the source's counit.
In the source the boundary algebra is the dual algebra $\mathcal A^*$ of a
pre-bialgebra $\mathcal A$, through the direct sum $\psi$ of the irreducible
representations of $\mathcal A^*$, and $T^{ik}=\psi(\phi_{ik})$ is the image
of the matrix coefficient $\phi_{ik}\colon x\mapsto\phi(x)_{ik}$ of the
injective physical representation $\phi$ of $\mathcal A$.  The counit of
$\mathcal A^*$ is $f\mapsto f(1)$ (lines~880 and~1419), and $\phi$ is a
representation of the unital algebra $\mathcal A$, so $\phi(1)=\Id$
(line~1421).  Hence the counit of $\mathcal B$ is
$\psi(f)\mapsto f(1)$, and $\psi(\phi_{ik})\mapsto\phi(1)_{ik}=\delta_{ik}$.

\section{What is missing}

A complete zipper fusion family records the blocks, the fusion tensors and
their left inverses of Ref.~\cite{Bultinck2017Anyons}, not the unital
algebra $\mathcal A$ or the unitality of $\phi$, and its fields do not force
a functional with~(\ref{eq:mrgscp22_counit_letters}) to exist.  Writing a
functional on $\mathcal B$ as $X\mapsto\sum_c\tr(Y_cX_c)$, such a functional
exists exactly when the identity operator on one site is a one-site operator
of the family,
\begin{align}
    \sum_c\tr\bigl(Y_cT_c^{ik}\bigr)
        &=\delta_{ik}
    \qquad\text{for some }Y\in\mathcal B.
    \label{eq:mrgscp22_unit_condition}
\end{align}
This can fail, and then $\Delta$ may have no counit at all.  Take one block
of bond dimension one, physical dimension two, letters
$T^{ik}=\delta_{i1}\delta_{k2}$, no fusion channel ($N^c_{ab}=0$, so the
fusion tensors are empty), and the left inverse that reads off the entry
$(i,k)=(1,2)$.  Every field holds: the block is injective since its only
nonzero letter spans $M_1(\C)$, and both zipper identities and the
reconstruction identity reduce to $\sum_jT^{ij}T^{jk}=0$.  The coproduct is
$\Delta=0$ on $\mathcal B=\C$, so the left side of~(\ref{eq:mrgscp22_counit})
vanishes for every $\epsilon$.

The formal pre-bialgebra structure on $\mathcal B$ therefore takes a
functional satisfying~(\ref{eq:mrgscp22_counit_letters}) as an argument.
With it, the statement is the source's statement that $\mathcal A^*$ is a
pre-bialgebra (line~1640), read on the boundary algebra.

\section{Elimination plan}

The condition~(\ref{eq:mrgscp22_unit_condition}) holds whenever the family
arises from the source's data, that is, from the irreducible representations
of the dual of a pre-bialgebra together with a unital injective physical
representation.  Formalizing that construction, from a pre-bialgebra to a
complete zipper fusion family, supplies the functional and removes the
argument.  Adding~(\ref{eq:mrgscp22_unit_condition}) as a field of the family
instead would add a hypothesis to the results stated for the family of
Ref.~\cite{Bultinck2017Anyons}, which does not assert it.

\bibliographystyle{alpha}
\bibliography{references}

\end{document}
